\documentclass[11pt,a4paper]{ctexart}
\usepackage[margin=25mm]{geometry}
\usepackage{amsmath,amssymb,amsthm,mathtools}
\usepackage{fancyhdr}
\usepackage{microtype}
\usepackage{enumitem}
\setlist{nosep}
\pagestyle{fancy}
\fancyhf{}
\fancyhead[L]{复分析 · 作业 1 / Complex Analysis · HW 1}
\fancyhead[R]{中英双语解答 / Bilingual Solutions}
\fancyfoot[C]{\thepage}
\setlength{\headheight}{14pt}
\setlength{\parindent}{0pt}
\setlength{\parskip}{5pt}
\newcommand{\dd}{\mathrm d}
\newcommand{\C}{\mathbb C}
\newcommand{\R}{\mathbb R}
\newcommand{\CN}{\textbf{中文解答.}\enspace}
\newcommand{\EN}{\textbf{English solution.}\enspace}
\begin{document}
\begin{center}
  {\LARGE\bfseries 复分析第 1 次作业解答}\\[5pt]
  {\large Complex Analysis Homework 1: Solutions}
\end{center}
\vspace{4pt}
\noindent\rule{\linewidth}{0.5pt}

\section*{1. 复线性与 Cauchy--Riemann 方程 / Complex linearity and the Cauchy--Riemann equations}
\textbf{题目 / Problem.} 设 $f=u+iv:\C\to\C$，$F=(u,v):\R^2\to\R^2$。
证明若 $f$ 全纯，则 $J_F$ 是复线性映射，且 Cauchy--Riemann 方程等价于
\[
 \begin{pmatrix}0&-1\\1&0\end{pmatrix}
 \begin{pmatrix}u_x&u_y\\v_x&v_y\end{pmatrix}
 =
 \begin{pmatrix}u_x&u_y\\v_x&v_y\end{pmatrix}
 \begin{pmatrix}0&-1\\1&0\end{pmatrix}.
\]
Let $f=u+iv:\C\to\C$ and $F=(u,v):\R^2\to\R^2$. Show that holomorphicity
of $f$ makes $J_F$ complex-linear, and that the Cauchy--Riemann equations
are equivalent to the displayed matrix identity.

\CN 令
\[
 A=J_F=\begin{pmatrix}u_x&u_y\\v_x&v_y\end{pmatrix},\qquad
 J=\begin{pmatrix}0&-1\\1&0\end{pmatrix}.
\]
在 $\R^2\cong\C$ 下，$J$ 正是“乘以 $i$”。实线性映射 $A$ 为复线性，当且仅当
$A(iw)=iA(w)$ 对每个 $w$ 成立，亦即 $AJ=JA$。直接相乘得到
\[
 JA=\begin{pmatrix}-v_x&-v_y\\u_x&u_y\end{pmatrix},\qquad
 AJ=\begin{pmatrix}u_y&-u_x\\v_y&-v_x\end{pmatrix}.
\]
两矩阵相等恰好等价于 $u_x=v_y$ 和 $u_y=-v_x$。若 $f$ 全纯，则实微分为
$h\mapsto f'(z)h$，所以 $J_F$ 复线性，以上等式也证明了所要求的矩阵表述。

\EN Under $\R^2\cong\C$, the matrix $J$ represents multiplication by $i$.
A real-linear map $A=J_F$ is complex-linear exactly when $A(iw)=iA(w)$ for every $w$,
or equivalently $AJ=JA$. The displayed products show that this commutation relation is
precisely $u_x=v_y$ and $u_y=-v_x$. If $f$ is holomorphic, its real differential is
$h\mapsto f'(z)h$, which is complex-linear. Thus its Jacobian commutes with $J$,
and the matrix identity is equivalent to the Cauchy--Riemann equations.

\section*{2. 极坐标形式 / Polar-coordinate form}
\textbf{题目 / Problem.} 证明 $f=u+iv$ 的极坐标 Cauchy--Riemann 方程为
$u_r=r^{-1}v_\theta$、$v_r=-r^{-1}u_\theta$。
Show that the Cauchy--Riemann equations in polar coordinates take the form
$u_r=r^{-1}v_\theta$ and $v_r=-r^{-1}u_\theta$.

\CN 以下取 $r>0$。由 $x=r\cos\theta$、$y=r\sin\theta$ 及链式法则，
\[
 u_r=u_x\cos\theta+u_y\sin\theta,
 \qquad \frac{v_\theta}{r}=-v_x\sin\theta+v_y\cos\theta.
\]
代入 $u_x=v_y$、$u_y=-v_x$，便得 $u_r=r^{-1}v_\theta$。同理，
\[
 v_r=v_x\cos\theta+v_y\sin\theta,
 \qquad -\frac{u_\theta}{r}=u_x\sin\theta-u_y\cos\theta,
\]
故 $v_r=-r^{-1}u_\theta$。由于从 $(r,\theta)$ 到 $(x,y)$ 的雅可比行列式为 $r\ne0$，
这些极坐标方程也反过来蕴含直角坐标的 Cauchy--Riemann 方程。

\EN For $r>0$, the chain rule gives the four identities displayed above.
Substituting $u_x=v_y$ and $u_y=-v_x$ yields
\[
 \boxed{\displaystyle u_r=\frac1r v_\theta,\qquad
 v_r=-\frac1r u_\theta.}
\]
Conversely, the polar change of variables has Jacobian determinant $r\ne0$;
reversing the same calculation recovers the Cartesian Cauchy--Riemann equations.
The formulas are not asserted at $r=0$, where polar coordinates are singular.

\section*{3. Laplace 算子与 Wirtinger 导数 / Laplacian and Wirtinger derivatives}
\textbf{题目 / Problem.} 证明 $\Delta=4\partial_z\partial_{\bar z}$。
Show that $\Delta=4\partial_z\partial_{\bar z}$.

\CN 对任意 $C^2$ 函数，定义
\[
 \partial_z=\frac12(\partial_x-i\partial_y),\qquad
 \partial_{\bar z}=\frac12(\partial_x+i\partial_y).
\]
由于混合偏导可交换，
\[
 4\partial_z\partial_{\bar z}
 =(\partial_x-i\partial_y)(\partial_x+i\partial_y)
 =\partial_x^2+\partial_y^2=\Delta.
\]
\EN For a $C^2$ function, expand the product of the two Wirtinger operators as above.
The mixed terms cancel because $\partial_x\partial_y=\partial_y\partial_x$,
leaving $4\partial_z\partial_{\bar z}=\Delta$.

\section*{4. 分式线性变换的分类 / Classification of fractional linear transformations}
\textbf{题目 / Problem.} 设标准化的分式线性变换
$T(z)=(az+b)/(cz+d)$ 有有限不动点 $z_1$，其乘子为 $\kappa=T'(z_1)$。
按题意，抛物型指 $\kappa=1$，双曲型指 $\kappa>0$ 且 $\kappa\ne1$，
椭圆型指 $|\kappa|=1$ 且 $\kappa\ne1$，其余为斜向型。
证明这四种情形依次等价于 $a+d=\pm2$、$a+d\in\R$ 且 $|a+d|>2$、
$a+d\in\R$ 且 $|a+d|<2$、以及 $a+d\notin\R$。
题目提示还要求讨论不动点数、第二个不动点的乘子及
$\kappa+\kappa^{-1}=(a+d)^2-2$。

Let $T(z)=(az+b)/(cz+d)$ be normalized and have a finite fixed point $z_1$.
Its multiplier is $\kappa=T'(z_1)$. The stated definitions are:
parabolic for $\kappa=1$; hyperbolic for positive real $\kappa\ne1$;
elliptic for $|\kappa|=1$, $\kappa\ne1$; and loxodromic otherwise.
Prove that these correspond, respectively, to trace $a+d=\pm2$,
real trace with $|a+d|>2$, real trace with $|a+d|<2$, and nonreal trace.
The hints concern the number of fixed points, the multiplier at a second
fixed point, and $\kappa+\kappa^{-1}=(a+d)^2-2$.

\CN “标准化”指 $ad-bc=1$。设 $t=a+d$、
$M=\begin{pmatrix}a&b\\c&d\end{pmatrix}\in\mathrm{SL}_2(\C)$。
若有限不动点 $z_1$ 的对应特征值是 $\lambda=cz_1+d$，则
\[
 M\binom{z_1}{1}=\lambda\binom{z_1}{1},\qquad
 \kappa=T'(z_1)=\frac{ad-bc}{(cz_1+d)^2}=\lambda^{-2}.
\]
另一特征值是 $\lambda^{-1}$，因此
\begin{equation}\label{eq:trace}
 \boxed{\displaystyle \kappa+\kappa^{-1}
 =\lambda^{-2}+\lambda^2=t^2-2.}
\end{equation}
若有两个不同不动点，它们的乘子互为倒数；第二个点为 $\infty$ 时，
在局部坐标 $w=1/z$ 中计算其乘子即可。若 $\kappa=1$，则 $t^2=4$，
矩阵有重特征值；除恒等变换外，$T$ 恰有一个不动点。

由式 \eqref{eq:trace} 可逐一判定：
\begin{align*}
 \kappa=1
 &\iff t=\pm2;\\
 \kappa>0,\ \kappa\ne1
 &\iff t\in\R,\ |t|>2;\\
 |\kappa|=1,\ \kappa\ne1
 &\iff t\in\R,\ |t|<2.
\end{align*}
例如第二行中 $\kappa+\kappa^{-1}>2$，故 $t^2>4$；反向则由
$X+X^{-1}=t^2-2>2$ 的两个根均为正实数可得。
第三行对应 $\kappa+\kappa^{-1}\in[-2,2)$，故 $t^2\in[0,4)$；
反向则两根在单位圆上。四类穷尽所有乘子，所以剩余的斜向型情形恰为 $t\notin\R$。

\EN Write $t=\operatorname{tr}M$ and normalize $M$ by $\det M=1$.
At a finite fixed point, $(z_1,1)^T$ is an eigenvector with eigenvalue
$\lambda=cz_1+d$, and $\kappa=T'(z_1)=\lambda^{-2}$.
The other eigenvalue is $\lambda^{-1}$, so equation \eqref{eq:trace} follows.
Two distinct fixed points therefore have reciprocal multipliers (use $w=1/z$
when one is at infinity). The alternatives now follow from the ranges of
$\kappa+\kappa^{-1}$: it equals $2$ for $\kappa=1$, lies in $(2,\infty)$
for positive real $\kappa\ne1$, and lies in $[-2,2)$ for unit-modulus
$\kappa\ne1$. Conversely, the quadratic equation
$X+X^{-1}=t^2-2$ gives the indicated type in each real-trace range.
All remaining cases have nonreal trace and are loxodromic.

\textbf{题目提示的例外 / Exception in the hint.}
恒等变换也满足 $\kappa=1$ 和 $t=\pm2$，但它的每个点都是不动点。
因此“$\kappa=1$ 当且仅当只有一个不动点”须排除恒等变换；
提示中第二个不动点处的 $f'(z_2)$ 应为 $T'(z_2)$。
The identity has multiplier $1$ and trace $\pm2$ but fixes every point;
the unique-fixed-point claim needs this exception. The derivative in the
second hint should be $T'(z_2)$.

\section*{5. Schwarzian 导数 / Schwarzian derivative}
\textbf{题目 / Problem.} 对全纯函数 $f$，其 Schwarzian 导数为
\[
 S(f)=\frac{f'''}{f'}-\frac32\left(\frac{f''}{f'}\right)^2.
\]
证明：(1) $S(f)\equiv0$ 当且仅当 $f$ 是分式线性变换；
(2) 对全纯 $g$，有 $S(g\circ f)=(S(g)\circ f)(f')^2+S(f)$，
特别地，若 $g$ 是分式线性变换，则 $S(g\circ f)=S(f)$。
提示：令 $N(f)=f''/f'$，则 $S(f)=N(f)'-\tfrac12N(f)^2$。

For a holomorphic $f$, define $S(f)=f'''/f'-\tfrac32(f''/f')^2$.
Show that (1) $S(f)\equiv0$ exactly for fractional linear maps and
(2) $S(g\circ f)=(S(g)\circ f)(f')^2+S(f)$ for holomorphic $g$;
in particular, postcomposition by a fractional linear map leaves
$S(f)$ unchanged. Hint: $N(f)=f''/f'$ satisfies
$S(f)=N(f)'-\tfrac12N(f)^2$.

\subsection*{(1) $S(f)\equiv0$ 与分式线性变换 / Vanishing Schwarzian and M\"obius maps}
\CN 在连通定义域上先假设 $f'\ne0$，令 $N=f''/f'$。直接求导得
\[
 S(f)=N'-\frac12N^2.
\]
若 $S(f)=0$，则 $N'=N^2/2$。如果 $N$ 在某点为零，微分方程解的唯一性
说明 $N\equiv0$，于是 $f(z)=Az+B$，其中 $A\ne0$。
否则 $N$ 处处非零，且
\[
 \left(\frac1N\right)'=-\frac12,
 \qquad N(z)=-\frac{2}{z-\zeta}
\]
（常数 $\zeta$ 不在定义域内）。积分两次得到
$f'(z)=C(z-\zeta)^{-2}$ 与
$f(z)=A(z-\zeta)^{-1}+B$，也是分式线性变换。
反过来，对于 $f(z)=(az+b)/(cz+d)$，
\[
 f'(z)=\frac{ad-bc}{(cz+d)^2},\qquad
 N(f)=-\frac{2c}{cz+d},\qquad
 N(f)'=\frac{2c^2}{(cz+d)^2}=\frac12N(f)^2,
\]
故 $S(f)=0$。

\EN On a connected domain where $f'\ne0$, put $N=f''/f'$.
Then $S(f)=N'-N^2/2$. If $S(f)=0$, the differential equation
$N'=N^2/2$ has either the zero solution, giving the affine map $Az+B$,
or a nowhere-zero solution $N=-2/(z-\zeta)$, giving
$f=A/(z-\zeta)+B$. Both are fractional linear maps.
Conversely, the displayed derivative calculation for
$f=(az+b)/(cz+d)$ gives $N'=N^2/2$, hence $S(f)=0$.
If critical points of $f$ were initially allowed, a zero Schwarzian on the
regular set actually rules them out: at a zero of $f'$ of order $m\ge1$,
$S(f)$ has leading term $-m(m+2)/(2(z-z_0)^2)$.

\subsection*{(2) 复合公式 / Composition law}
\CN 在相关导数非零处，设 $A=(N(g)\circ f)f'$、$B=N(f)$。
链式法则给出 $N(g\circ f)=A+B$，且
\[
 \begin{aligned}
 S(g\circ f)
 &=A'+B'-\tfrac12(A+B)^2\\
 &=\bigl(N(g)'-\tfrac12N(g)^2\bigr)\circ f\,(f')^2
   +\bigl(B'-\tfrac12B^2\bigr)
   +(N(g)\circ f)\bigl(f''-f'B\bigr)\\
 &=(S(g)\circ f)(f')^2+S(f),
 \end{aligned}
\]
因为 $f'B=f''$。若 $g$ 是分式线性变换，由第 (1) 问 $S(g)=0$，
遂有 $S(g\circ f)=S(f)$。

\EN The chain rule gives $N(g\circ f)=(N(g)\circ f)f'+N(f)$.
Expanding $N(g\circ f)'-\tfrac12N(g\circ f)^2$ as above,
the cross term cancels because $f'N(f)=f''$. Thus
\[
 \boxed{\displaystyle S(g\circ f)=(S(g)\circ f)(f')^2+S(f).}
\]
For a fractional linear $g$, part (1) gives $S(g)=0$, so
$S(g\circ f)=S(f)$.

\end{document}
